%% ma: L12-B08-0137
%% lop: 12
%% chuong: 2
%% bai: 8
%% dang: 12.8.47
%% loai: TN4
%% muc: TH
%% dap_an: "D"
%% nguon: "CĐ2.3.pdf giáo viên gửi 27/10/2026 (M07, Lê Minh Khôi – Nguyễn Anh Khôi, Chương II Chủ đề 3 Biểu thức tọa độ của các phép toán vectơ), Câu 52"
%% kien_thuc: "Bài 8 (khoảng cách hai điểm), lớp 10 (định lí côsin)"
%% trang_thai: chinh-thuc
%% ghi_chu: "Giáo viên đã duyệt 27/10/2026"
\begin{ex}
Trong không gian $Oxyz$, cho ba điểm $A(1;2;0)$, $B(1;0;-1)$, $C(0;-1;2)$. Tam giác $ABC$ là
\choice{tam giác vuông đỉnh $A$}{tam giác đều}{tam giác cân đỉnh $A$}{\True tam giác có ba góc nhọn}
\loigiai{$AB^2=5$, $BC^2=11$, $CA^2=14$: ba cạnh khác nhau.

Mỗi cạnh bình phương nhỏ hơn tổng bình phương hai cạnh kia ($14<16$) nên tam giác có ba góc nhọn.}
\end{ex}
